\section{线性方程与最小二乘方法的比较}

前面的算法都围绕线性关系展开，但返回的对象并不相同。选择方法前，应先确定任务要求满足方程、最小化残差，还是在多个最优解中选出长度最小的一个；随后再检查矩阵结构和可用的数据访问方式。

\subsection{求解目标与方法适用性}

相容方程要求原残差为零。若零空间非平凡，这仍不能唯一决定未知量：从一个解加上任何零空间向量，输出都不变。最小范数要求解与零空间正交，是额外的选解条件。

不相容问题的原残差不能消失。欧氏最小二乘要求正规残差为零；如果又要求最小范数，还须使解位于行空间。因而“残差下降”“达到最小二乘解”和“得到最小范数解”是三种不同的判断，不能用同一个停止量代替。

\subsection{矩阵条件}

表中“矩阵作用”指能够计算矩阵与向量的乘积，不必保存整个矩阵。迭代法的条件指本章已经建立保证的情形；一般矩阵上的表现不能由方法名称保证。

\begin{center}
\small
\begin{tabular}{p{0.24\textwidth}p{0.65\textwidth}}
\toprule
方法 & 适用条件与主要输出\\
\midrule
LU & 可逆方阵；选主元消元后求精确方程的解\\
Cholesky & 对称正定方阵；分解可供多个右端复用\\
块LDL$^\top$ & 对称可逆矩阵，可含不定性；需要对称选主元\\
Schur补 & 已知分块且被消去的块可逆；先求约化方程\\
Woodbury & 已有可复用的基矩阵求解器，更新秩较低\\
Richardson & 对称正定时可按谱范围选步长；只用矩阵作用\\
Jacobi & 对角非零；严格行对角占优保证收敛，正定本身不够\\
Gauss--Seidel & 严格行对角占优或对称正定；逐行原位更新\\
SOR、SSOR & 对称正定且$0<\omega<2$；顺序扫掠或双向扫掠\\
CG & 对称正定；短递推，最小化能量范数误差\\
GMRES & 一般非奇异方阵；最小化原残差，需保存搜索基\\
MINRES & 对称非奇异方阵，可不定；最小化原残差\\
BiCG & 一般非对称方阵；需要矩阵及转置作用，可能失效\\
BiCGSTAB & 一般非对称方阵；无需转置作用，仍须检查失效\\
\bottomrule
\end{tabular}
\end{center}

矩形问题还需要比较求解目标。下面几种方法直接以最小二乘或行投影为出发点，不必先把原矩阵改成方阵。

\begin{center}
\small
\begin{tabular}{p{0.24\textwidth}p{0.65\textwidth}}
\toprule
方法 & 适用条件与主要输出\\
\midrule
正规方程法 & 满列秩最小二乘；形成正规矩阵会平方条件数\\
QR & 满列秩最小二乘；正交变换后求三角方程\\
SVD & 包括秩亏情形；给出秩、最小残差与最小范数代表\\
LSQR、LSMR & 矩形矩阵与转置作用；分别侧重原残差和正规残差\\
Kaczmarz & 相容系统，逐行访问；零初值选出最小范数解\\
扩展Kaczmarz & 允许不相容，需行列访问及辅助向量；概率保证另论\\
\bottomrule
\end{tabular}
\end{center}

左乘可逆矩阵保持精确方程的解集，却一般改变不相容问题中残差的权重。右换元保持可达到的输出，但欧氏最小范数随坐标尺度改变。把问题转换成某算法能够接收的形状以后，还要检查求解目标是否仍是原来的目标。

\subsection{计算开销与收敛速度}

稠密$n$阶LU、Cholesky、QR分解的工作量都为$O(n^3)$，因子的存储为$O(n^2)$；Cholesky利用对称性减少常数。分解完成后，每个新右端通常只需$O(n^2)$的三角求解。对$m\ge d$的稠密矩形矩阵，QR的主要工作量为$O(md^2)$。稀疏情形不能直接把这些估计中的维数换成非零元数，因为消元可能产生填充。

若矩阵有$s$个非零元，一次矩阵作用通常为$O(s)$。CG与BiCGSTAB等短递推保存常数个工作向量；完整GMRES的第$k$步还需对$k$个旧向量正交化，增加$O(kn)$工作和$O(kn)$存储。LSQR与LSMR每步使用一次矩阵作用和一次转置作用，工作向量总长度与$m+d$同阶。重正交化会改变这一存储估计。

因此总成本是每轮代价乘以达到要求所需的轮数，还要加上预条件构造和应用成本。只比较迭代次数，或只比较一次矩阵作用，都不足以判断哪个方法更省。

\subsection{数值稳定性}

正交变换避免了换基中的长度放大，因而QR常比显式形成正规方程更可靠；它仍不能消除原矩阵的小奇异值。SVD揭示这些方向，却也不能从带噪数据中恢复本来就不确定的分量。截断和惩罚会改变问题，应与单纯改善计算过程区别开来。

迭代结束时重新计算原残差。最小二乘还检查正规残差，并结合矩阵尺度解释容差。递推残差很小而真实残差较大，说明舍入误差已影响递推；相邻解几乎不变则可能只是停滞。若数据本身只精确到一定尺度，过度降低计算残差未必改善最终解。

\subsection{典型场景中的取舍}

对中等规模的稠密可逆方阵，可以从带选主元LU开始；若已知对称正定，改用Cholesky。多个右端共用矩阵时，复用分解通常比重复启动迭代更合适。已有分解而矩阵只发生低秩修改时，再检查Woodbury的小系统是否可逆、是否值得复用。

对大型稀疏对称正定系统，可以使用预条件CG；对称不定时考虑保持对称结构的MINRES。一般非对称系统中，GMRES用较多存储换取明确的残差最小化，BiCGSTAB用短递推减少存储，但需接受残差波动及失效检查。

对稠密矩形最小二乘，满列秩时可用QR；需要识别秩或最小范数时使用SVD。大型矩形问题若提供矩阵与转置作用，可用LSQR或LSMR；只方便逐行读取且系统相容时，行投影才直接符合数据访问方式。算法的选择由目标、结构与成本共同决定。
